\section{Critical-Cycle Characterization}

Let $V$ contain one vertex for each orbit of payoff coordinates $(i,a)$ under $G$.  For every player $i$, opponent profile $a_{-i}$, and ordered pair of distinct own actions $(b_i,a_i)$, insert an edge
\[
e:[i,(b_i,a_{-i})]\longrightarrow[i,(a_i,a_{-i})]
\]
with label
\[
c_e=u_i(a_i,a_{-i})-u_i(b_i,a_{-i}).
\]
The reverse ordered comparison is also present before orbit identification.  Parallel edges are allowed; for a fixed ordered endpoint pair, only the smallest label matters for the one-sided upper-bound system.  Self-loops must be retained because symmetry identification can collapse a nonzero comparison onto one vertex.

An invariant surrogate assigns one potential $x_v$ to every vertex, with
$\widehat u_i(a)=x_{[i,a]}$.  Its error on edge $e$ is
\[
|x_{\mathrm{head}(e)}-x_{\mathrm{tail}(e)}-c_e|.
\]

\begin{theorem}[Maximum-mean-cycle formula]
\label{thm:supp-cycle}
Let $\mathcal C$ be the directed cycles of the orbit graph, including self-loops.  With value zero when $\mathcal C$ is empty,
\begin{equation}
\delta_G(\Gamma)
=
\max_{C\in\mathcal C}
\left(-\frac{1}{|C|}\sum_{e\in C}c_e\right).
\label{eq:supp-cycle}
\end{equation}
At this value, difference-constraint potentials construct a nearest invariant surrogate, while a maximizing cycle certifies that no smaller error is possible.
\end{theorem}
\begin{proof}
A potential has strategic error at most $t$ exactly when
\begin{equation}
-t
\le
x_{\mathrm{head}(e)}-x_{\mathrm{tail}(e)}-c_e
\le t
\label{eq:supp-two-sided}
\end{equation}
for every ordered comparison.  The upper inequality is
\[
x_{\mathrm{head}(e)}-x_{\mathrm{tail}(e)}\le c_e+t.
\]
The lower inequality is the upper inequality for the reverse comparison.  Thus \eqref{eq:supp-two-sided} is equivalent to a directed system of difference constraints with edge length $c_e+t$.

Such a system is feasible iff every directed cycle has nonnegative total length.  Necessity follows by summing the inequalities around a cycle, where the potentials telescope.  For sufficiency, add a zero-length source edge to every vertex and use shortest-path distances; when no negative cycle exists, these distances are finite and satisfy every constraint.  Therefore $t$ is feasible iff
\[
0\le\sum_{e\in C}(c_e+t)
=
\sum_{e\in C}c_e+|C|t
\qquad\forall C,
\]
or equivalently
\[
t\ge-rac1{|C|}\sum_{e\in C}c_e
\qquad\forall C.
\]
Minimizing over $t\ge0$ proves \eqref{eq:supp-cycle}.  At the critical value, shortest-path distances for lengths $c_e+\delta_G$ give the required orbit potentials.
\end{proof}

\subsection{Certified Witness Extraction}

Karp's dynamic program computes the maximum cycle mean of weights $w_e=-c_e$ in polynomial time \cite{karp1978cycle}.  A predecessor chain in the value table need not itself be a critical cycle, so the implementation uses a separate tight-edge argument.

Let $\lambda=\delta_G$ and define reduced weights $w'_e=w_e-\lambda$.  Every directed cycle has reduced total weight at most zero, and at least one critical cycle has total weight zero whenever the graph contains a positive optimum.  Add a source with zero-weight edges to every vertex and define $h_v$ as the maximum reduced weight of a directed walk from the source to $v$.  Because there is no positive cycle, these maxima are finite and satisfy
\begin{equation}
h_{\mathrm{head}(e)}
\ge
h_{\mathrm{tail}(e)}+w'_e.
\label{eq:supp-maxplus}
\end{equation}
Call an edge tight when equality holds.  The slack in \eqref{eq:supp-maxplus} is nonnegative.  Around a zero-weight critical cycle, the potential differences telescope and the edge weights sum to zero, so the slacks sum to zero.  Every edge on that critical cycle is therefore tight.  A directed-cycle search in the tight-edge subgraph returns a checkable witness whose mean is exactly $\lambda$.  The code then verifies both closure and witness mean numerically.

\subsection{Worked Two-Orbit Example}

Suppose two payoff-coordinate orbits $U,V$ receive comparisons
$U\to V$ with label $1$ and $V\to U$ with label $-3$.  Writing
$d=x_V-x_U$, the two errors are
\[
|d-1|
\qquad\text{and}\qquad
|3-d|.
\]
Their minimax value is one, attained at $d=2$.  The two-edge cycle has mean label $(1-3)/2=-1$, and \eqref{eq:supp-cycle} gives the same defect.  The certificate identifies the source of inconsistency: after the symmetry coordinates are merged into only $U$ and $V$, one context asks their potential difference to be near $1$, while another asks it to be near $3$.

\subsection{Cheap Screening Bounds}

For any generator set $S$,
\begin{equation}
\frac12\max_{s\in S}\devdist(\Gamma,s\Gamma)
\le
\delta_G(\Gamma)
\le
\devdist(\Gamma,\mathcal P_G\Gamma).
\label{eq:supp-screening}
\end{equation}
The upper bound evaluates the feasible group-averaged surrogate.  For the lower bound, any invariant $\widehat\Gamma$ satisfies
\[
\devdist(\Gamma,s\Gamma)
\le
\devdist(\Gamma,\widehat\Gamma)
+
\devdist(s\widehat\Gamma,s\Gamma)
=
2\devdist(\Gamma,\widehat\Gamma).
\]
Thus generator discrepancies and orbit averaging can screen a hierarchy before running the sharp projection.  A critical cycle can then guide refinement by identifying which orbit identifications or payoff estimates dominate the remaining defect.
